\begin{frame}{자주 묻는 질문(15.2)}
  \small
  \begin{itemize}
    \item \textbf{Q: 틱택토(완전탐색)와 님(MCTS)은 왜 다른
      방법?} --- 틱택토는 549,946개로 완전탐색이 가능해
      minimax로 ``100\% 정확한'' 답. 님은 \textbf{사실}
      완전탐색이 가능하지만, \textbf{일부러} MCTS로 풀어
      ``완전탐색 없이 표본추출만으로도 최적 전략을 찾아낼 수
      있다''를 보여주기 위한 예제 --- 실전에서 MCTS를 쓰는
      이유는 바둑처럼 완전탐색이 \textbf{아예 불가능한}
      문제이기 때문
    \item \textbf{Q: MCTS는 지도학습? 강화학습?} --- \textbf{엄밀히
      둘 다 아님} --- \textbf{계획(planning)} 알고리즘.
      Week 4 가치반복처럼 모델(규칙)이 있다는 전제하에
      ``지금 무엇을 할지''를 계산, Week 9--11처럼 경사하강으로
      파라미터를 갱신하는 게 \textbf{아님}. AlphaZero가
      강한 이유 = 정확히 이 계획(MCTS)과 학습(신경망
      지도학습)을 \textbf{루프로 엮었기 때문}
    \item \textbf{Q: ``가장 많이 방문된 자식'' vs ``승률 최대
      자식''?} --- 표준은 \textbf{방문 최대(most visits)}.
      방문이 적을 때 승률 추정치는 \textbf{노이져스}(20번 중
      3승 vs 100번 중 30승은 승률 같아도 신뢰도 다름),
      방문 횟수는 UCT가 ``여기를 더 봐야겠다''고 판단한
      \textbf{양} 자체가 견고. 시뮬레이션이 충분히 많으면
      대개 같은 수를 가리키고, 예산이 끝났을 때 승률 격차가
      확실히 큰 \textbf{예외 상황}에서만 승률 최대 변형
  \end{itemize}
\end{frame}
